% Generated draft: requires mathematical and source review.
\documentclass[11pt]{article}
\usepackage[margin=25mm]{geometry}
\usepackage{fontspec}
\setmainfont{Latin Modern Roman}
\usepackage{amsmath,amssymb,mathtools,unicode-math}
\setmathfont{Latin Modern Math}
\usepackage{enumitem,needspace,xurl,hyperref}
\setlength{\parindent}{0pt}
\setlength{\parskip}{5pt}
\setlength{\emergencystretch}{3em}
\newcommand{\N}{\mathbb{N}}
\newcommand{\Z}{\mathbb{Z}}
\newcommand{\Q}{\mathbb{Q}}
\newcommand{\R}{\mathbb{R}}
\newcommand{\C}{\mathbb{C}}
\newcommand{\F}{\mathbb{F}}
\newcommand{\A}{\mathbb{A}}
\newcommand{\PP}{\mathbb{P}}
\newcommand{\E}{\mathbb{E}}
\newcommand{\eps}{\varepsilon}
\newcommand{\dd}{\,\mathrm{d}}
\begin{document}
\section*{UnsolvedMath ID: 30000570}
\textbf{Source problem code:} catalogue entry
\subsection{Definitions and mathematical statement}
In a usual Banach-space setting, let \(X\) be a separable infinite-dimensional Banach space over \(\mathbb K\in\{\R,\C\}\), and let \(T:X\to X\) be a bounded linear operator. A vector \(x\) is supercyclic for \(T\) when its projective orbit is dense:
\[
 \overline{\{\lambda T^n x:\lambda\in\mathbb K,\ n\in\Z_{\geq0}\}}=X.
\]
A supercyclic subspace conventionally means a closed infinite-dimensional linear subspace \(L\subseteq X\) such that every \(0\ne x\in L\) is supercyclic for \(T\). This is stronger than existence of a single supercyclic vector, and it does not impose \(T\)-invariance on \(L\). The source does not specify the scalar field, Banach versus more general topological-vector-space setting, or the closed/infinite-dimensional convention, so these definitions are stated as the common interpretation rather than imposed source hypotheses.

Let \(\mathsf P_{\mathrm{spec}}(T)\) denote the particular spectral condition associated with the Supercyclicity Criterion in the report. That condition and the version of the criterion are not printed in the supplied record. The question is to identify natural operator classes \(\mathcal C\) for which
\[
 \forall T\in\mathcal C,\qquad
 [\exists L\text{ a supercyclic subspace for }T]
 \ \Longrightarrow\ \mathsf P_{\mathrm{spec}}(T).
\]
This asks necessity of the specified spectral condition within such a class. It does not ask for sufficiency alone, does not assert necessity for all bounded operators, and does not require the full Supercyclicity Criterion to be necessary merely because one associated spectral condition is being tested.

The source does not say whether the condition uses the spectrum, an essential spectrum, a particular spectral radius or a geometry of spectral components. No guessed spectral intersection or radius inequality is substituted. ``Natural class'' also has no stated closure, structural or nonvacuity requirements; simply defining a class by the desired implication would not identify the research target. A definite statement needs the exact spectral predicate and the intended concrete operator families.

The literature triage checked 2026-08-21 reports no authoritative resolution of the extracted claim located in its targeted search. The exact-statement label says the sentence was retained byte-for-byte. Those labels are supplied metadata, not an independently verified current theorem or a source for the omitted criterion.
\subsection{Short English statement}
Find meaningful operator families in which having a closed infinite-dimensional subspace of supercyclic vectors forces the source's spectral condition. Necessity is the requested direction. The condition itself, operator setting and intended meaning of natural class are missing, so no particular spectrum or inequality is guessed.
\subsection{Sources}
\begin{itemize}
\item Local UnsolvedMath registry, permanent ID 30000570, source code catalogue entry, ``Spectral Criteria for Supercyclic Subspaces,'' including the unchanged original necessity question and exact-statement verification.
\item \emph{Mini-Workshop: Hypercyclicity and Linear Chaos (2007)}, Oberwolfach Reports, pp.\ 2227--2276, DOI 10.4171/owr/2006/37: \url{https://doi.org/10.4171/owr/2006/37}. This is also the sole supplied literature reference. The citation uses 2007, while the DOI and proposed-year metadata use 2006.
\item Literature triage checked 2026-08-21; local record updated 2026-09-16. No external source has been visited or independently checked here.
\end{itemize}
\end{document}
